\scanpage{177}

Pour prouver (ii), on est ramené à prouver que t\textsuperscript{t} élément de $\mathrm{Prépic}'(X'/S')$ dont les images sont les \unsure{mêmes|$\text{m}^{\mathrm{es}}$}, provient d'un élément de $\mathrm{Prépic}'(X/S)$. À cause de l'unicité, on voit que c'est local sur $S$, on peut alors supposer qu'on a une section de $X$ sur $S$.

On peut alors identifier les \deleted{objets}\inserted{élémts} de $\mathrm{Prépic}\,(X/S)$ à des faisceaux inversibles \inserted{(déts à \uline{isom unique} près)} sur $X$ à section marquée sur $S$, ce qui permet d'appliquer la théorie de la descente \uline{fid plate}.

\uline{Corollaire 2} \unsure{Bien|On} \unsure{combiner|combine} avec Prop 1. (ii) pour donner des \uline{faisceaux} sur $X$.

\uline{Corollaire \uline{3}} Supposons que la relation $H^0_{X/S}(\underline O_X)=\underline O_S$ reste vraie pour tte extension de la base (i.e., $X$ \unsure{coh.|est} plat sur $S$ en dim 0, et $\underline H^0_{X/S}(\underline O_X)=\underline O_S$), \inserted{et $X/S$ a localement une section} Alors le foncteur $\text{\uuline{Prépic}}'_{X/S}$ est compatible avec [= exact pour] \deleted{\unsure{une}} la descente fidèlement plate quasi-compacte.

% archive sheet number: 176
